% !TEX root = ../main.tex
\section{Internal Representations and What the Argument Actually Needs}
\label{sec:internal}

\subsection{The premise at issue}

Section~\ref{sec:structural-argument} identified the premise that carries the
argument from a claim about the output operator to a claim about the system:
that there is no prompt-invariant state \emph{governing} assertion for the
prompt-form to elicit. I flagged there that the premise is weaker than the
denial of any prompt-invariant state whatever. This section explains why the
weaker premise is the right one, and why the stronger one should not be
attempted.

The stronger premise is contested by a substantial body of work. Unsupervised
methods identify directions in activation space that behave like truth-values
across a range of contents \citep{burns2023}. Classifiers trained on internal
activations can distinguish assertions the model will make that are true from
those that are false, and can do so before the assertion is produced
\citep{azaria2023}. If a direction in activation space tracks truth across
paraphrases, then there is something prompt-invariant in there, and a flat denial
is simply false.
% TODO(cite): the geometry-of-truth line of work belongs here. I could not verify
% venue and year for Marks & Tegmark; check before citing.

I do not want to deny it, for two reasons. The first is that the evidence points
the other way and a philosophical argument that requires an empirical claim to be
false is a bad bet. The second is better: the evidence, read carefully, supports
the premise I need.

\subsection{Representation without governance}
\label{sec:dissociation}

Consider what the strongest result in this literature actually reports. A
classifier on internal activations can tell that the model is about to assert
something false. The model asserts it anyway. That is not evidence that assertion
expresses an internal attitude. It is evidence of a \emph{dissociation}: the
internal state registers one thing and the assertion mechanism does another.

This is exactly the structure the account of Section~\ref{sec:subdoxastic}
predicts. A state that carries information about truth but does not govern what
the system asserts is a state that fails the constitutive-aim condition on
belief. Belief aims at truth in the sense that recognising a candidate as false
is, for a believer, sufficient reason not to hold it and not to assert it. A
system in which the recognition and the assertion come apart---in which the first
can be read off the activations while the second proceeds regardless---is a system
where nothing plays the functional role that makes closure and consistency norms
rather than accidents.

So the premise to state is not that the model has no prompt-invariant states. It
is that no prompt-invariant state is \emph{accountable} for what the model
asserts. There is no mechanism whose function is to make assertion answer to
whatever the activations represent. Put in the vocabulary of
Section~\ref{sec:grounding}: the missing thing is precisely condition (2) of
Definition~\ref{def:grounding}, generalised from evidence to internal state.
Grounding is a demand for accountability of assertion to something; the finding
that a truth-correlated representation exists and is not that something is the
sharpest available statement of what these systems lack.

There is an irony worth noting. The research programme that looks most like
evidence for machine belief---probing for internal truth representations---is
also the programme that has documented most precisely the gap between what the
system represents and what it says. If the representations are real, the
dissociation is real, and the dissociation is the phenomenon.

\subsection{What would have to be shown}
\label{sec:probing-limits}

The weaker premise is not free either, and honesty requires noting how far the
probing literature is from settling matters in any direction. The empirical and
conceptual obstacles to reading probe results as belief attributions are
substantial: probes may track features correlated with truth rather than truth,
may be artefacts of the training distribution used to fit them, and may not
generalise across content domains in the way a belief-representation claim
requires \citep{levinstein2024}. Standards for what would count as a genuine
belief representation---rather than a decodable correlate---have been set out
explicitly \citep{herrmann2024}, and they are demanding.

I take the state of play to be this. Whether these systems have
belief-\emph{representations} is open, and I do not need to close it. Whether
those representations govern assertion is not open in the same way: the observed
dissociation is direct evidence that they do not. My thesis needs only the second
claim, and it is the second claim that the practical consequences follow from. A
system whose assertions are not accountable to its own best internal estimate of
truth is a system whose assertions must be made accountable to something
external---an identified, applicable evidential base---if they are to be relied
upon at all. That is the argument of
Sections~\ref{sec:measurement} and~\ref{sec:warrant}, and it survives whatever the
probing literature eventually concludes about representation.

One consequence for the objections that follow. If the premise is the weaker one,
then the interpretationist of Section~\ref{sec:interpretationism} cannot
be answered by denying that there is anything inside; the answer must be about
predictive projectibility, which is where I take it. And the human-parallel
objection of Section~\ref{sec:human-parallel} becomes sharper rather than
weaker, because the human case is precisely one in which recognising a claim as
false does govern whether one asserts it.
